% =========================================================
% IMPOSTAZIONI GENERALI
% =========================================================

\usepackage[utf8]{inputenc}
\usepackage[T1]{fontenc}
\usepackage[italian]{babel}

\usepackage{lmodern}
\usepackage{microtype}

% =========================================================
% MATEMATICA
% =========================================================

\usepackage{amsmath}
\usepackage{amssymb}
\usepackage{amsthm}
\usepackage{mathtools}

% Insiemi numerici
\newcommand{\N}{\mathbb{N}}
\newcommand{\Z}{\mathbb{Z}}
\newcommand{\Q}{\mathbb{Q}}
\newcommand{\R}{\mathbb{R}}
\newcommand{\C}{\mathbb{C}}

% Valore assoluto e norma
\newcommand{\abs}[1]{\left|#1\right|}
\newcommand{\norm}[1]{\left\lVert#1\right\rVert}

% =========================================================
% GRAFICA
% =========================================================

\usepackage{graphicx}
\usepackage{float}

% Disegni matematici
\usepackage{tikz}

% Grafici di funzioni
\usepackage{pgfplots}
\pgfplotsset{compat=1.18}

% Riquadri
\usepackage[most]{tcolorbox}

% Librerie TikZ utili
\usetikzlibrary{
    arrows.meta,
    calc,
    intersections,
    patterns,
    positioning,
    decorations.pathreplacing
}

% =========================================================
% TEOREMI E DEFINIZIONI
% =========================================================

\theoremstyle{definition}
\newtheorem{definizione}{Definizione}[section]
\newtheorem{esempio}{Esempio}[section]
\newtheorem{esercizio}{Esercizio}[section]

\theoremstyle{plain}
\newtheorem{teorema}{Teorema}[section]
\newtheorem{proposizione}{Proposizione}[section]
\newtheorem{lemma}{Lemma}[section]
\newtheorem{corollario}{Corollario}[section]

% =========================================================
% NOTE
% =========================================================

\newtcolorbox{nota}{
    breakable,
    title={Nota bene},
    fonttitle=\bfseries,
    colback=blue!4,
    colframe=blue!55!black,
    boxrule=0.8pt,
    arc=2mm,
    left=3mm,
    right=3mm,
    top=2mm,
    bottom=2mm
}

% =========================================================
% IMPAGINAZIONE
% =========================================================

\usepackage[
    a4paper,
    margin=2.5cm
]{geometry}

\usepackage{enumitem}

% =========================================================
% LINK
% =========================================================

\usepackage{hyperref}

\hypersetup{
    colorlinks=true,
    linkcolor=blue,
    urlcolor=blue,
    citecolor=blue
}

% =========================================================
% FINE LEZIONE
% =========================================================

\newcommand{\finelezione}{
    \vspace{1cm}
    \begin{center}
        $\diamond\qquad\diamond\qquad\diamond$

        \textit{Fine della lezione}
    \end{center}
}

% =========================================================
% APPROFONDIMENTI PERSONALI
% =========================================================

\newtcolorbox{approfondimento}[1][]{
    breakable,
    title={Approfondimento personale},
    fonttitle=\bfseries,
    colback=gray!5,
    colframe=gray!60,
    boxrule=0.6pt,
    arc=2mm,
    #1
}

% =========================================================
% FINE APPROFONDIMENTO
% =========================================================

\newcommand{\fineapprofondimento}{
    \vspace{1cm}
    \begin{center}
        $\diamond\qquad\diamond\qquad\diamond$

        \textit{Fine dell'approfondimento}
    \end{center}
}
